\textbf{结论名称：}圆幂定理统一形式

\textbf{适用条件：}已知定圆 \(O\) 半径 \(R\) 及平面内一点 \(P\)，记 \(d = OP\)。过点 \(P\) 作任意直线交圆于 \(A, B\)（切、交均可；重合时视为切线情形）。结论对点 \(P\) 在圆内、圆上、圆外均成立。

\textbf{变量说明：}圆心 \(O\)，圆半径 \(R > 0\)；点 \(P\) 为平面内任意一点，\(d = OP \geq 0\)；过 \(P\) 的直线与圆交于 \(A\) 与 \(B\)（若不交则只讨论虚交点，此处仅考虑实交或相切）；\(PA\)、\(PB\) 表示两交点（或切点）到 \(P\) 的线段长度（非有向）。

\textbf{核心公式：}
\[
\boxed{PA \cdot PB = |d^{2} - R^{2}|}
\]

\textbf{使用场景：}统一处理相交弦定理、割线定理、切割线定理；在竞赛或模拟题中快速求与端点、切点相关的线段乘积，或由乘积反求点到圆心的距离，避免分情况讨论。

\textbf{简单例子：}设圆 \(O\) 半径 \(R=5\)，点 \(P\) 在圆内且 \(OP=3\)。过 \(P\) 任作弦 \(AB\)，求 \(PA \cdot PB\)。解：由公式得 \(PA \cdot PB = |3^{2} - 5^{2}| = |9 - 25| = 16\)。可验证：若取过 \(P\) 的直径，其两段分别为 \(2\) 与 \(8\)，乘积亦为 \(16\)。